% !TEX program = xelatex
\documentclass[UTF8,11pt,twoside,openany]{ctexbook}
\usepackage[paperwidth=176mm,paperheight=250mm,inner=19.5mm,outer=19.5mm,top=20mm,bottom=20mm,headheight=16pt,headsep=6mm,footskip=10mm]{geometry}
\usepackage{amsmath,amssymb,bm,array,longtable,booktabs,fancyhdr,hyperref}
\hypersetup{hidelinks,pdftitle={物理竞赛二轮讲义：十编合订本},pdfauthor={}}
\linespread{1.18}\setlength{\parindent}{2em}\setlength{\parskip}{3pt}\setlength{\emergencystretch}{2em}\raggedbottom
\ctexset{chapter={format=\LARGE\bfseries\centering,beforeskip=16pt,afterskip=22pt},section={format=\large\bfseries}}
\pagestyle{fancy}\fancyhf{}\fancyfoot[C]{\small 总-\thepage}\fancyhead[C]{\small 物理竞赛二轮讲义}\renewcommand{\headrulewidth}{0pt}
\fancypagestyle{plain}{\fancyhf{}\fancyfoot[C]{\small 总-\thepage}\renewcommand{\headrulewidth}{0pt}}
\begin{document}\frontmatter
\begin{titlepage}\centering\vspace*{35mm}
{\Huge\bfseries 物理竞赛二轮讲义\par}
\vspace{14mm}{\Large 十编合订本\par}
\vspace{20mm}{\large 数学方法与微积分\par 力学与电磁学\par 热学、近代物理与光学\par}
\vfill{\large 2026年11月修订\par}
\end{titlepage}
\chapter*{前言}\addcontentsline{toc}{chapter}{前言}
二轮课程在已有定律的基础上讨论条件更复杂的模型。本讲义先说明研究对象、约束和近似条件，再建立方程、展开计算，并考察边界或初态改变后解法的变化。守恒关系始终与所选系统对应。

全书分十编，共四十五讲。第一编讲复数、几何、线性代数与AI辅助计算；第二、三编讨论运动、静力与动力学综合；第四编补充物理问题中使用的微积分方法；第五至七编依次讨论静电静磁、电磁感应和电路；第八至十编为热学、近代物理和光学。各编可以单独使用，合订本保留分编目录和局部章、题、图号，便于与分册对照。

这里的“讲”按内容划分，模型较多的专题可分次讨论。原课程的知识、模型和题设条件构成各编主线，扩充内容补足推导与变式。题解遇到已证通式时给出引用，新的建模步骤和条件转换仍逐项说明。

第一编的向量、矩阵和微分方法是后续各编共同使用的工具。阅读第三编的小振动和第七编的网络谱方法时，可回查第一编第三讲的本征矢与正交展开；处理第四编的受迫振动和第七编的暂态电路时，可对照二阶线性方程的初值解。电磁学中的边界条件、热学中的系统边界，以及光学中的波矢与能流，各有具体物理含义，不能仅因方程形式相近就互相替换。

例题与习题均附解答。题设图确定几何、方向和变量，解答图接续受力分解、过程变化或边界构造，并与方程中的符号对应。可选题目标为“选做”，理论中的“选读”表示扩展阅读；公式的适用条件均在正文中说明。

页码采用“编号—本编页码”，例如“7-80”表示第七编第80页。跨编引用应同时给出编名和讲次；“第七编，习题2.4”与该编分册中的同一道题对应。全书总目录列出四十五讲及主要节题，主题与模型索引用于查找分散在不同编中的相关内容。电子版目录的页码和书签均可跳转，纸面页码与分册保持一致。

AI辅助计算用于实现已经明确的物理模型。方程、单位、初始条件和终止条件应先写清，再用解析极限、误差估计与守恒关系核对输出。数值曲线与推导不符时，应回到模型和算法查明原因；计算结果本身不能代替物理解释。
\chapter*{全书目录}\addcontentsline{toc}{chapter}{全书目录}
本目录中的页码均为“编号—本编页码”。各编前置说明与详细目录使用本编罗马页码，正文使用阿拉伯页码。
\section*{第一编\quad 大学物理中的数学方法}
\begin{longtable}{@{}p{.84\textwidth}r@{}}
\textbf{前言}&\href{course-page:1:2}{1-ii}\\[2pt]
\textbf{第一讲　复数、相位与复变函数}&\href{course-page:1:6}{1-1}\\[2pt]
\hspace{1em}从二维几何到复数运算&\href{course-page:1:6}{1-1}\\[2pt]
\hspace{1em}极形式、旋转与开方&\href{course-page:1:8}{1-3}\\[2pt]
\hspace{1em}欧拉公式与复指数&\href{course-page:1:10}{1-5}\\[2pt]
\hspace{1em}同频振动与相量&\href{course-page:1:11}{1-6}\\[2pt]
\hspace{1em}复指数方法的运算边界&\href{course-page:1:14}{1-9}\\[2pt]
\hspace{1em}受迫阻尼振动：从瞬态到稳态&\href{course-page:1:16}{1-11}\\[2pt]
\hspace{1em}阻尼模型的完整比较&\href{course-page:1:19}{1-14}\\[2pt]
\hspace{1em}近似的尺度：共振附近该保留什么&\href{course-page:1:22}{1-17}\\[2pt]
\hspace{1em}交流电路中的复阻抗&\href{course-page:1:27}{1-22}\\[2pt]
\hspace{1em}并联电路、复功率与相位三角形&\href{course-page:1:29}{1-24}\\[2pt]
\hspace{1em}傅里叶展开与多频线性响应&\href{course-page:1:31}{1-26}\\[2pt]
\hspace{1em}复可导、解析性与调和函数&\href{course-page:1:34}{1-29}\\[2pt]
\hspace{1em}复平面上的轨迹与映射&\href{course-page:1:36}{1-31}\\[2pt]
\hspace{1em}复积分、奇点与留数（选读）&\href{course-page:1:39}{1-34}\\[2pt]
\hspace{1em}复积分的路径选择与估计（选读）&\href{course-page:1:42}{1-37}\\[2pt]
\hspace{1em}习题与解答&\href{course-page:1:45}{1-40}\\[2pt]
\textbf{第二讲　向量、曲线与空间几何}&\href{course-page:1:65}{1-60}\\[2pt]
\hspace{1em}向量的分量与几何意义&\href{course-page:1:65}{1-60}\\[2pt]
\hspace{1em}直线、平面与空间距离&\href{course-page:1:68}{1-63}\\[2pt]
\hspace{1em}圆锥曲线的几何定义&\href{course-page:1:69}{1-64}\\[2pt]
\hspace{1em}抛体轨迹、可达区域与包络&\href{course-page:1:71}{1-66}\\[2pt]
\hspace{1em}极坐标运动学与中心力&\href{course-page:1:73}{1-68}\\[2pt]
\hspace{1em}曲线的切向、法向与曲率&\href{course-page:1:76}{1-71}\\[2pt]
\hspace{1em}旋转坐标与惯性力&\href{course-page:1:79}{1-74}\\[2pt]
\hspace{1em}从约束选坐标：抛物线轨道上的小珠&\href{course-page:1:81}{1-76}\\[2pt]
\hspace{1em}面积元、体积元与雅可比&\href{course-page:1:86}{1-81}\\[2pt]
\hspace{1em}立体角、通量与对称性&\href{course-page:1:89}{1-84}\\[2pt]
\hspace{1em}连续分布的积分：圆环与圆盘&\href{course-page:1:89}{1-84}\\[2pt]
\hspace{1em}曲线积分与曲面积分的对象&\href{course-page:1:91}{1-86}\\[2pt]
\hspace{1em}梯度、方向导数与等值面&\href{course-page:1:94}{1-89}\\[2pt]
\hspace{1em}散度、旋度与局部测量&\href{course-page:1:96}{1-91}\\[2pt]
\hspace{1em}格林、散度与斯托克斯定理&\href{course-page:1:97}{1-92}\\[2pt]
\hspace{1em}曲线坐标中的微分算子（选读）&\href{course-page:1:101}{1-96}\\[2pt]
\hspace{1em}复势与二维无源场（选读）&\href{course-page:1:104}{1-99}\\[2pt]
\hspace{1em}习题与解答&\href{course-page:1:106}{1-101}\\[2pt]
\textbf{第三讲　线性代数与物理系统的结构}&\href{course-page:1:126}{1-121}\\[2pt]
\hspace{1em}线性组合、独立性与基&\href{course-page:1:126}{1-121}\\[2pt]
\hspace{1em}矩阵、线性变换与方程组&\href{course-page:1:127}{1-122}\\[2pt]
\hspace{1em}行列式、逆与几何缩放&\href{course-page:1:129}{1-124}\\[2pt]
\hspace{1em}正交、投影与换基&\href{course-page:1:130}{1-125}\\[2pt]
\hspace{1em}非正交坐标与度量矩阵&\href{course-page:1:132}{1-127}\\[2pt]
\hspace{1em}特征值、特征向量与对角化&\href{course-page:1:135}{1-130}\\[2pt]
\hspace{1em}线性变换的几何与特征方向&\href{course-page:1:137}{1-132}\\[2pt]
\hspace{1em}二次型与平衡点稳定性&\href{course-page:1:137}{1-132}\\[2pt]
\hspace{1em}两个耦合振子的正常模&\href{course-page:1:140}{1-135}\\[2pt]
\hspace{1em}质量不同、零频模与质量归一化&\href{course-page:1:143}{1-138}\\[2pt]
\hspace{1em}不同质量双振子的完整求解&\href{course-page:1:145}{1-140}\\[2pt]
\hspace{1em}弱耦合、失谐与能量交换&\href{course-page:1:146}{1-141}\\[2pt]
\hspace{1em}多振子链与边界条件&\href{course-page:1:152}{1-147}\\[2pt]
\hspace{1em}转动惯量张量与主轴&\href{course-page:1:153}{1-148}\\[2pt]
\hspace{1em}矩阵指数与线性微分方程（选读）&\href{course-page:1:155}{1-150}\\[2pt]
\hspace{1em}强迫耦合系统与反共振&\href{course-page:1:158}{1-153}\\[2pt]
\hspace{1em}习题与解答&\href{course-page:1:161}{1-156}\\[2pt]
\textbf{第四讲　AI 辅助物理建模与计算实验}&\href{course-page:1:183}{1-178}\\[2pt]
\hspace{1em}从一个问题得到可运行的实验&\href{course-page:1:183}{1-178}\\[2pt]
\hspace{1em}简谐振子的数值误差&\href{course-page:1:184}{1-179}\\[2pt]
\hspace{1em}实验一：单摆的小角近似能用到哪里&\href{course-page:1:187}{1-182}\\[2pt]
\hspace{1em}实验二：共振峰为什么不在同一处&\href{course-page:1:189}{1-184}\\[2pt]
\hspace{1em}实验三：弱耦合与失谐&\href{course-page:1:190}{1-185}\\[2pt]
\hspace{1em}实验四：有阻力抛体与最佳发射角&\href{course-page:1:192}{1-187}\\[2pt]
\hspace{1em}实验五：从数据识别模型遗漏&\href{course-page:1:193}{1-188}\\[2pt]
\hspace{1em}选读：采样让两个信号看起来相同&\href{course-page:1:196}{1-191}\\[2pt]
\hspace{1em}事件定位与计算步长&\href{course-page:1:196}{1-191}\\[2pt]
\hspace{1em}参数反演中的可识别性&\href{course-page:1:199}{1-194}\\[2pt]
\hspace{1em}静摩擦切换与分段运动&\href{course-page:1:201}{1-196}\\[2pt]
\hspace{1em}事件终点的灵敏度与参数优化&\href{course-page:1:203}{1-198}\\[2pt]
\hspace{1em}随机抽样的误差与物理测度&\href{course-page:1:206}{1-201}\\[2pt]
\hspace{1em}稀有事件的加权抽样&\href{course-page:1:210}{1-205}\\[2pt]
\hspace{1em}习题与解答&\href{course-page:1:213}{1-208}\\[2pt]
\textbf{综合习题}&\href{course-page:1:231}{1-226}\\[2pt]
\textbf{习题提示}&\href{course-page:1:236}{1-231}\\[2pt]
\hspace{1em}第一章&\href{course-page:1:236}{1-231}\\[2pt]
\hspace{1em}第二章&\href{course-page:1:236}{1-231}\\[2pt]
\hspace{1em}第三章&\href{course-page:1:237}{1-232}\\[2pt]
\hspace{1em}第四章&\href{course-page:1:238}{1-233}\\[2pt]
\textbf{符号与计算约定}&\href{course-page:1:239}{1-234}\\[2pt]
\end{longtable}
\section*{第二编\quad 运动和静力进阶}
\begin{longtable}{@{}p{.84\textwidth}r@{}}
\textbf{第一讲　微元法进阶}&\href{course-page:2:3}{2-1}\\[2pt]
\hspace{1em}微元选取与变量变换&\href{course-page:2:3}{2-1}\\[2pt]
\hspace{1em}角分布、立体角与靶面密度&\href{course-page:2:12}{2-10}\\[2pt]
\hspace{1em}连续受力与碰撞散射&\href{course-page:2:20}{2-18}\\[2pt]
\hspace{1em}习题与解答&\href{course-page:2:30}{2-28}\\[2pt]
\textbf{第二讲　运动关联进阶}&\href{course-page:2:47}{2-45}\\[2pt]
\hspace{1em}绳杆约束与刚体关联&\href{course-page:2:47}{2-45}\\[2pt]
\hspace{1em}接触与滚动关联&\href{course-page:2:55}{2-53}\\[2pt]
\hspace{1em}复杂机构与转动参考系&\href{course-page:2:61}{2-59}\\[2pt]
\hspace{1em}接触点迁移与交点运动&\href{course-page:2:69}{2-67}\\[2pt]
\hspace{1em}习题与解答&\href{course-page:2:74}{2-72}\\[2pt]
\textbf{第三讲　曲线与三维运动}&\href{course-page:2:92}{2-90}\\[2pt]
\hspace{1em}曲率与自然坐标&\href{course-page:2:92}{2-90}\\[2pt]
\hspace{1em}刚体平面运动&\href{course-page:2:98}{2-96}\\[2pt]
\hspace{1em}空间运动&\href{course-page:2:104}{2-102}\\[2pt]
\hspace{1em}动坐标与相对运动&\href{course-page:2:113}{2-111}\\[2pt]
\hspace{1em}空间滚动与定点转动&\href{course-page:2:125}{2-123}\\[2pt]
\hspace{1em}习题与解答&\href{course-page:2:130}{2-128}\\[2pt]
\textbf{第四讲　相图的应用}&\href{course-page:2:147}{2-145}\\[2pt]
\hspace{1em}能量与相轨线&\href{course-page:2:147}{2-145}\\[2pt]
\hspace{1em}平衡点与分界轨线&\href{course-page:2:156}{2-154}\\[2pt]
\hspace{1em}阻尼及综合应用&\href{course-page:2:165}{2-163}\\[2pt]
\hspace{1em}守恒关系与参数平面的允许区域&\href{course-page:2:179}{2-177}\\[2pt]
\hspace{1em}习题与解答&\href{course-page:2:188}{2-186}\\[2pt]
\textbf{第五讲　自由度和能量解法}&\href{course-page:2:210}{2-208}\\[2pt]
\hspace{1em}自由度与坐标选择&\href{course-page:2:210}{2-208}\\[2pt]
\hspace{1em}虚功与静力平衡&\href{course-page:2:218}{2-216}\\[2pt]
\hspace{1em}约束系统的能量解法&\href{course-page:2:227}{2-225}\\[2pt]
\hspace{1em}柔性系统与杆架的虚功&\href{course-page:2:238}{2-236}\\[2pt]
\hspace{1em}电磁系统中的虚功与外源功&\href{course-page:2:245}{2-243}\\[2pt]
\hspace{1em}习题与解答&\href{course-page:2:252}{2-250}\\[2pt]
\end{longtable}
\section*{第三编\quad 动力学综合}
\begin{longtable}{@{}p{.84\textwidth}r@{}}
\textbf{第一讲　能量解法问题}&\href{course-page:3:4}{3-1}\\[2pt]
\hspace{1em}能量积分、坐标加速度与约束反力&\href{course-page:3:4}{3-1}\\[2pt]
\hspace{1em}纯滚动系统的动能与接触条件&\href{course-page:3:16}{3-13}\\[2pt]
\hspace{1em}多自由度、循环坐标与整体反冲&\href{course-page:3:21}{3-18}\\[2pt]
\hspace{1em}滑动摩擦与滚动状态的转换&\href{course-page:3:26}{3-23}\\[2pt]
\hspace{1em}习题与解答&\href{course-page:3:30}{3-27}\\[2pt]
\textbf{第二讲　约束临界问题}&\href{course-page:3:51}{3-48}\\[2pt]
\hspace{1em}单侧约束与临界分离&\href{course-page:3:51}{3-48}\\[2pt]
\hspace{1em}刚体支承、滑动与内部受力&\href{course-page:3:57}{3-54}\\[2pt]
\hspace{1em}约束极值、虚功与二阶变化&\href{course-page:3:63}{3-60}\\[2pt]
\hspace{1em}曲线约束与最速下降&\href{course-page:3:71}{3-68}\\[2pt]
\hspace{1em}移动约束的反力与分离条件&\href{course-page:3:75}{3-72}\\[2pt]
\hspace{1em}习题与解答&\href{course-page:3:78}{3-75}\\[2pt]
\textbf{第三讲　质心运动问题}&\href{course-page:3:93}{3-90}\\[2pt]
\hspace{1em}固定物质系统的质心与整体反冲&\href{course-page:3:93}{3-90}\\[2pt]
\hspace{1em}质量输运、动量与能量收支&\href{course-page:3:101}{3-98}\\[2pt]
\hspace{1em}固有加速度与相对论推进&\href{course-page:3:112}{3-109}\\[2pt]
\hspace{1em}流束偏转中的动量通量与功率&\href{course-page:3:118}{3-115}\\[2pt]
\hspace{1em}习题与解答&\href{course-page:3:120}{3-117}\\[2pt]
\textbf{第四讲　简正模问题}&\href{course-page:3:135}{3-132}\\[2pt]
\hspace{1em}平衡构形、二次能量与固有频率&\href{course-page:3:135}{3-132}\\[2pt]
\hspace{1em}振型分解、简正坐标与初始条件&\href{course-page:3:149}{3-146}\\[2pt]
\hspace{1em}拍振、失谐与外力响应&\href{course-page:3:159}{3-156}\\[2pt]
\hspace{1em}零频模态与线性近似的边界&\href{course-page:3:165}{3-162}\\[2pt]
\hspace{1em}习题与解答&\href{course-page:3:170}{3-167}\\[2pt]
\textbf{第五讲　刚体碰撞进阶}&\href{course-page:3:180}{3-177}\\[2pt]
\hspace{1em}碰撞守恒量、冲量与有效质量&\href{course-page:3:180}{3-177}\\[2pt]
\hspace{1em}斜碰、速度空间与允许末态&\href{course-page:3:189}{3-186}\\[2pt]
\hspace{1em}多约束、摩擦与连续碰撞&\href{course-page:3:196}{3-193}\\[2pt]
\hspace{1em}习题与解答&\href{course-page:3:213}{3-210}\\[2pt]
\textbf{第六讲　刚体运动进阶}&\href{course-page:3:222}{3-219}\\[2pt]
\hspace{1em}惯量张量、主轴与角动量&\href{course-page:3:222}{3-219}\\[2pt]
\hspace{1em}欧拉角、进动与刚体方程&\href{course-page:3:231}{3-228}\\[2pt]
\hspace{1em}平面刚体、运动点取矩与滚动&\href{course-page:3:241}{3-238}\\[2pt]
\hspace{1em}加速与转动参考系中的约束运动&\href{course-page:3:250}{3-247}\\[2pt]
\hspace{1em}习题与解答&\href{course-page:3:271}{3-268}\\[2pt]
\textbf{第七讲　连续体问题}&\href{course-page:3:304}{3-301}\\[2pt]
\hspace{1em}连续体的微元、弹性与边界功&\href{course-page:3:304}{3-301}\\[2pt]
\hspace{1em}行波、反射与动态边界&\href{course-page:3:311}{3-308}\\[2pt]
\hspace{1em}驻波、模态正交与初始条件&\href{course-page:3:324}{3-321}\\[2pt]
\hspace{1em}梁的弯曲、曲率与端部条件&\href{course-page:3:329}{3-326}\\[2pt]
\hspace{1em}有限夹层中的反射、透射与阻抗匹配&\href{course-page:3:331}{3-328}\\[2pt]
\hspace{1em}弯曲波的色散与能量输运&\href{course-page:3:333}{3-330}\\[2pt]
\hspace{1em}分布质量对悬臂梁低频振动的影响&\href{course-page:3:336}{3-333}\\[2pt]
\hspace{1em}习题与解答&\href{course-page:3:338}{3-335}\\[2pt]
\textbf{第八讲　天体运动进阶}&\href{course-page:3:348}{3-345}\\[2pt]
\hspace{1em}二体运动、有效势与径向分支&\href{course-page:3:348}{3-345}\\[2pt]
\hspace{1em}由初态确定轨道的形状与方向&\href{course-page:3:356}{3-353}\\[2pt]
\hspace{1em}面积定律、飞行时间与会合相位&\href{course-page:3:364}{3-361}\\[2pt]
\hspace{1em}瞬时变轨、碰撞截面与散射&\href{course-page:3:370}{3-367}\\[2pt]
\hspace{1em}轨道微扰、进动与潮汐&\href{course-page:3:377}{3-374}\\[2pt]
\hspace{1em}习题与解答&\href{course-page:3:386}{3-383}\\[2pt]
\end{longtable}
\section*{第四编\quad 微积分进阶}
\begin{longtable}{@{}p{.84\textwidth}r@{}}
\textbf{第一讲　一元微积分与渐近展开}&\href{course-page:4:3}{4-1}\\[2pt]
\hspace{1em}导数、微分与物理过程&\href{course-page:4:3}{4-1}\\[2pt]
\hspace{1em}泰勒展开与保留阶数&\href{course-page:4:12}{4-10}\\[2pt]
\hspace{1em}积分、换元与边界&\href{course-page:4:18}{4-16}\\[2pt]
\hspace{1em}展开、积分与小参数的配合&\href{course-page:4:26}{4-24}\\[2pt]
\hspace{1em}习题与解答&\href{course-page:4:32}{4-30}\\[2pt]
\textbf{第二讲　微分方程与振动}&\href{course-page:4:58}{4-56}\\[2pt]
\hspace{1em}由过程建立方程&\href{course-page:4:58}{4-56}\\[2pt]
\hspace{1em}线性方程与初始条件&\href{course-page:4:61}{4-59}\\[2pt]
\hspace{1em}非齐次方程与外力响应&\href{course-page:4:66}{4-64}\\[2pt]
\hspace{1em}约束运动中的降阶与终止条件&\href{course-page:4:75}{4-73}\\[2pt]
\hspace{1em}解的结构与非线性方程的局部行为&\href{course-page:4:78}{4-76}\\[2pt]
\hspace{1em}习题与解答&\href{course-page:4:81}{4-79}\\[2pt]
\textbf{第三讲　多元微分与场论}&\href{course-page:4:112}{4-110}\\[2pt]
\hspace{1em}偏导数、全微分与复合函数&\href{course-page:4:112}{4-110}\\[2pt]
\hspace{1em}无约束极值与约束极值&\href{course-page:4:116}{4-114}\\[2pt]
\hspace{1em}梯度、散度与旋度&\href{course-page:4:135}{4-133}\\[2pt]
\hspace{1em}二次变分与曲线坐标的几何&\href{course-page:4:141}{4-139}\\[2pt]
\hspace{1em}习题与解答&\href{course-page:4:143}{4-141}\\[2pt]
\textbf{第四讲　多重积分与积分变换}&\href{course-page:4:166}{4-164}\\[2pt]
\hspace{1em}曲线积分与曲面积分&\href{course-page:4:166}{4-164}\\[2pt]
\hspace{1em}多重积分与区域划分&\href{course-page:4:173}{4-171}\\[2pt]
\hspace{1em}把边界积分化为内部积分&\href{course-page:4:183}{4-181}\\[2pt]
\hspace{1em}高斯积分、Gamma与Beta函数&\href{course-page:4:188}{4-186}\\[2pt]
\hspace{1em}傅里叶与拉普拉斯变换&\href{course-page:4:192}{4-190}\\[2pt]
\hspace{1em}变换中的守恒关系与收敛条件&\href{course-page:4:199}{4-197}\\[2pt]
\hspace{1em}延伸讨论：离散采样与频谱&\href{course-page:4:201}{4-199}\\[2pt]
\hspace{1em}习题与解答&\href{course-page:4:202}{4-200}\\[2pt]
\end{longtable}
\section*{第五编\quad 静电静磁进阶}
\begin{longtable}{@{}p{.84\textwidth}r@{}}
\textbf{第一讲　场线、等势面与多极场}&\href{course-page:5:3}{5-1}\\[2pt]
\hspace{1em}源、场与电势&\href{course-page:5:3}{5-1}\\[2pt]
\hspace{1em}场的微分与曲线坐标&\href{course-page:5:13}{5-11}\\[2pt]
\hspace{1em}静电关系在完整场方程中的位置&\href{course-page:5:23}{5-21}\\[2pt]
\hspace{1em}由场能到导体表面的电压力&\href{course-page:5:24}{5-22}\\[2pt]
\hspace{1em}习题与解答&\href{course-page:5:27}{5-25}\\[2pt]
\textbf{第二讲　唯一性、局部展开与边值问题}&\href{course-page:5:46}{5-44}\\[2pt]
\hspace{1em}边界条件与唯一性&\href{course-page:5:46}{5-44}\\[2pt]
\hspace{1em}局部展开与约束下的稳定性&\href{course-page:5:49}{5-47}\\[2pt]
\hspace{1em}多极展开及原点的选择&\href{course-page:5:54}{5-52}\\[2pt]
\hspace{1em}二维边值问题与分离变量&\href{course-page:5:61}{5-59}\\[2pt]
\hspace{1em}边界几何造成的场增强与角模衰减&\href{course-page:5:66}{5-64}\\[2pt]
\hspace{1em}边界误差怎样传入内部区域&\href{course-page:5:70}{5-68}\\[2pt]
\hspace{1em}习题与解答&\href{course-page:5:71}{5-69}\\[2pt]
\textbf{第三讲　电像法、电容与导电类比}&\href{course-page:5:88}{5-86}\\[2pt]
\hspace{1em}电像法的构造与物理区域&\href{course-page:5:88}{5-86}\\[2pt]
\hspace{1em}双球的无限像链&\href{course-page:5:96}{5-94}\\[2pt]
\hspace{1em}两族圆与偏心圆柱电容&\href{course-page:5:103}{5-101}\\[2pt]
\hspace{1em}旋转椭球与圆盘的电容&\href{course-page:5:105}{5-103}\\[2pt]
\hspace{1em}稳恒导电与绝缘空穴&\href{course-page:5:109}{5-107}\\[2pt]
\hspace{1em}浮置导体的消元与机械功&\href{course-page:5:114}{5-112}\\[2pt]
\hspace{1em}习题与解答&\href{course-page:5:116}{5-114}\\[2pt]
\textbf{第四讲　电介质与静磁场}&\href{course-page:5:132}{5-130}\\[2pt]
\hspace{1em}极化强度、电位移与束缚电荷&\href{course-page:5:132}{5-130}\\[2pt]
\hspace{1em}界面边界条件&\href{course-page:5:135}{5-133}\\[2pt]
\hspace{1em}电容中的分区与能量&\href{course-page:5:140}{5-138}\\[2pt]
\hspace{1em}介质边界中的电像与多极解&\href{course-page:5:146}{5-144}\\[2pt]
\hspace{1em}静磁场、磁化与边界&\href{course-page:5:153}{5-151}\\[2pt]
\hspace{1em}介质壳层与带气隙磁路的完整接续&\href{course-page:5:157}{5-155}\\[2pt]
\hspace{1em}习题与解答&\href{course-page:5:160}{5-158}\\[2pt]
\end{longtable}
\section*{第六编\quad 电磁感应进阶}
\begin{longtable}{@{}p{.84\textwidth}r@{}}
\textbf{第一讲　电磁感应与场的表示}&\href{course-page:6:4}{6-1}\\[2pt]
\hspace{1em}回路、电动势与磁通的变化&\href{course-page:6:4}{6-1}\\[2pt]
\hspace{1em}局部场方程与微分运算&\href{course-page:6:6}{6-3}\\[2pt]
\hspace{1em}电势、磁矢势与规范&\href{course-page:6:7}{6-4}\\[2pt]
\hspace{1em}带电粒子的能量与正则动量&\href{course-page:6:10}{6-7}\\[2pt]
\hspace{1em}感生电场的空间分布&\href{course-page:6:12}{6-9}\\[2pt]
\hspace{1em}刚性导体的运动与无阻回路&\href{course-page:6:20}{6-17}\\[2pt]
\hspace{1em}感应电压的测量路径与滑动接触&\href{course-page:6:26}{6-23}\\[2pt]
\hspace{1em}习题与解答&\href{course-page:6:30}{6-27}\\[2pt]
\textbf{第二讲　自感、互感与磁路}&\href{course-page:6:48}{6-45}\\[2pt]
\hspace{1em}磁通、磁链与电感的定义&\href{course-page:6:48}{6-45}\\[2pt]
\hspace{1em}磁能、磁共能与外源功&\href{course-page:6:49}{6-46}\\[2pt]
\hspace{1em}互感、互易与耦合强度&\href{course-page:6:56}{6-53}\\[2pt]
\hspace{1em}耦合电感的连接与退化情形&\href{course-page:6:58}{6-55}\\[2pt]
\hspace{1em}磁路、气隙与磁阻网络&\href{course-page:6:63}{6-60}\\[2pt]
\hspace{1em}变压器中的共同磁通与分支磁通&\href{course-page:6:67}{6-64}\\[2pt]
\hspace{1em}非线性磁链与有损副边的反作用&\href{course-page:6:71}{6-68}\\[2pt]
\hspace{1em}习题与解答&\href{course-page:6:75}{6-72}\\[2pt]
\textbf{第三讲　电磁耦合与超导}&\href{course-page:6:93}{6-90}\\[2pt]
\hspace{1em}回旋运动与磁矩的绝热不变量&\href{course-page:6:93}{6-90}\\[2pt]
\hspace{1em}导电、弛豫与移动磁场&\href{course-page:6:95}{6-92}\\[2pt]
\hspace{1em}几何对称性与互感的分解&\href{course-page:6:100}{6-97}\\[2pt]
\hspace{1em}从理想导电到超导磁响应&\href{course-page:6:104}{6-101}\\[2pt]
\hspace{1em}导轨中的电磁耦合方程&\href{course-page:6:110}{6-107}\\[2pt]
\hspace{1em}有限响应时间与可逆储能&\href{course-page:6:123}{6-120}\\[2pt]
\hspace{1em}习题与解答&\href{course-page:6:127}{6-124}\\[2pt]
\textbf{第四讲　电磁场综合应用}&\href{course-page:6:147}{6-144}\\[2pt]
\hspace{1em}由互感求力与磁偶极近似&\href{course-page:6:147}{6-144}\\[2pt]
\hspace{1em}位移电流、传播场与能量流&\href{course-page:6:153}{6-150}\\[2pt]
\hspace{1em}电场、磁场与参考系&\href{course-page:6:156}{6-153}\\[2pt]
\hspace{1em}磁能的形变导数与电磁应力&\href{course-page:6:158}{6-155}\\[2pt]
\hspace{1em}磁场边界、转动与受约束轨道&\href{course-page:6:160}{6-157}\\[2pt]
\hspace{1em}场的边界与有限响应&\href{course-page:6:167}{6-164}\\[2pt]
\hspace{1em}习题与解答&\href{course-page:6:174}{6-171}\\[2pt]
\end{longtable}
\section*{第七编\quad 电路进阶}
\begin{longtable}{@{}p{.84\textwidth}r@{}}
\textbf{第一讲　微观导电与稳恒电路}&\href{course-page:7:4}{7-1}\\[2pt]
\hspace{1em}从载流子运动到局部欧姆定律&\href{course-page:7:4}{7-1}\\[2pt]
\hspace{1em}连续性、场方程与界面&\href{course-page:7:6}{7-3}\\[2pt]
\hspace{1em}非均匀材料中的电荷与电流&\href{course-page:7:8}{7-5}\\[2pt]
\hspace{1em}导电与静电的边值类比&\href{course-page:7:17}{7-14}\\[2pt]
\hspace{1em}从连续导体到节点方程&\href{course-page:7:25}{7-22}\\[2pt]
\hspace{1em}真实电源的等效与内循环电流&\href{course-page:7:29}{7-26}\\[2pt]
\hspace{1em}功率关系与端口描述&\href{course-page:7:31}{7-28}\\[2pt]
\hspace{1em}耗散原理与内部环流&\href{course-page:7:32}{7-29}\\[2pt]
\hspace{1em}横向电场与张量形式的欧姆定律&\href{course-page:7:35}{7-32}\\[2pt]
\hspace{1em}习题与解答&\href{course-page:7:37}{7-34}\\[2pt]
\textbf{第二讲　对称性、互易与网络变换}&\href{course-page:7:47}{7-44}\\[2pt]
\hspace{1em}唯一性与等势节点的归并&\href{course-page:7:47}{7-44}\\[2pt]
\hspace{1em}三端网络的星三角变换&\href{course-page:7:50}{7-47}\\[2pt]
\hspace{1em}对称性之外的桥式网络&\href{course-page:7:53}{7-50}\\[2pt]
\hspace{1em}线性叠加与正多面体&\href{course-page:7:56}{7-53}\\[2pt]
\hspace{1em}相邻端口电阻之和&\href{course-page:7:62}{7-59}\\[2pt]
\hspace{1em}互易定理与多端口&\href{course-page:7:67}{7-64}\\[2pt]
\hspace{1em}电阻变化的灵敏度&\href{course-page:7:71}{7-68}\\[2pt]
\hspace{1em}轮形网络与无限递推&\href{course-page:7:75}{7-72}\\[2pt]
\hspace{1em}逐层内接三角形的自洽电阻&\href{course-page:7:80}{7-77}\\[2pt]
\hspace{1em}有限矩形网络的边界谱与端口电阻&\href{course-page:7:82}{7-79}\\[2pt]
\hspace{1em}习题与解答&\href{course-page:7:85}{7-82}\\[2pt]
\textbf{第三讲　差分方程、无限网络与分形电路}&\href{course-page:7:93}{7-90}\\[2pt]
\hspace{1em}有限梯形与端部条件&\href{course-page:7:93}{7-90}\\[2pt]
\hspace{1em}双无限网络的共模与差模&\href{course-page:7:97}{7-94}\\[2pt]
\hspace{1em}横档与轨道不等阻时的衰减尺度&\href{course-page:7:101}{7-98}\\[2pt]
\hspace{1em}周期边界与环形双轨&\href{course-page:7:108}{7-105}\\[2pt]
\hspace{1em}四角正方形分形的三端等效&\href{course-page:7:112}{7-109}\\[2pt]
\hspace{1em}有限层分形与无限极限&\href{course-page:7:117}{7-114}\\[2pt]
\hspace{1em}从离散梯形到有泄漏的连续电阻带&\href{course-page:7:124}{7-121}\\[2pt]
\hspace{1em}非均匀泄漏带的幂律与局部近似&\href{course-page:7:127}{7-124}\\[2pt]
\hspace{1em}习题与解答&\href{course-page:7:130}{7-127}\\[2pt]
\textbf{第四讲　多轨网络与非稳态电路}&\href{course-page:7:138}{7-135}\\[2pt]
\hspace{1em}三轨电阻带的横向模态&\href{course-page:7:138}{7-135}\\[2pt]
\hspace{1em}模态方法中的端口选择&\href{course-page:7:140}{7-137}\\[2pt]
\hspace{1em}圆柱网络与离散傅里叶分解&\href{course-page:7:141}{7-138}\\[2pt]
\hspace{1em}电容电路的独立变量与初态&\href{course-page:7:148}{7-145}\\[2pt]
\hspace{1em}两只电容并不总是一阶&\href{course-page:7:154}{7-151}\\[2pt]
\hspace{1em}串联RLC的二阶暂态&\href{course-page:7:156}{7-153}\\[2pt]
\hspace{1em}二阶电路的一般初态与能量&\href{course-page:7:161}{7-158}\\[2pt]
\hspace{1em}正弦驱动与自由暂态&\href{course-page:7:166}{7-163}\\[2pt]
\hspace{1em}二极管与非线性负载线&\href{course-page:7:167}{7-164}\\[2pt]
\hspace{1em}非线性工作点的解析形式与局部响应&\href{course-page:7:170}{7-167}\\[2pt]
\hspace{1em}纯RC网络的衰减谱与非单调节点电压&\href{course-page:7:173}{7-170}\\[2pt]
\hspace{1em}习题与解答&\href{course-page:7:175}{7-172}\\[2pt]
\end{longtable}
\section*{第八编\quad 热学进阶}
\begin{longtable}{@{}p{.84\textwidth}r@{}}
\textbf{第一讲　分子运动的统计描述}&\href{course-page:8:4}{8-1}\\[2pt]
\hspace{1em}从状态概率到速度分布&\href{course-page:8:4}{8-1}\\[2pt]
\hspace{1em}速率、方向与抽样方式&\href{course-page:8:9}{8-6}\\[2pt]
\hspace{1em}相对运动、碰撞与自由程&\href{course-page:8:15}{8-12}\\[2pt]
\hspace{1em}活塞、分子流与壁面通量&\href{course-page:8:18}{8-15}\\[2pt]
\hspace{1em}碰撞抽样与运动边界的进一步应用&\href{course-page:8:27}{8-24}\\[2pt]
\hspace{1em}自由飞行、方向随机化与扩散尺度&\href{course-page:8:37}{8-34}\\[2pt]
\hspace{1em}习题与解答&\href{course-page:8:41}{8-38}\\[2pt]
\textbf{第二讲　泄流、辐射与热力学关系}&\href{course-page:8:48}{8-45}\\[2pt]
\hspace{1em}小孔泄流中的能量选择&\href{course-page:8:48}{8-45}\\[2pt]
\hspace{1em}光子气体的压强与辐射&\href{course-page:8:55}{8-52}\\[2pt]
\hspace{1em}热力学势与偏导关系&\href{course-page:8:57}{8-54}\\[2pt]
\hspace{1em}热容、状态方程与过程约束的综合应用&\href{course-page:8:69}{8-66}\\[2pt]
\hspace{1em}有限换热速度下的气体压缩&\href{course-page:8:78}{8-75}\\[2pt]
\hspace{1em}习题与解答&\href{course-page:8:82}{8-79}\\[2pt]
\textbf{第三讲　第二定律、熵与有限热源}&\href{course-page:8:90}{8-87}\\[2pt]
\hspace{1em}热机、卡诺循环与效率上界&\href{course-page:8:90}{8-87}\\[2pt]
\hspace{1em}熵的定义与不可逆过程&\href{course-page:8:93}{8-90}\\[2pt]
\hspace{1em}混合气体与有限物体的熵变&\href{course-page:8:98}{8-95}\\[2pt]
\hspace{1em}分批接触、换热与温度转移&\href{course-page:8:107}{8-104}\\[2pt]
\hspace{1em}过程设计与可用功&\href{course-page:8:114}{8-111}\\[2pt]
\hspace{1em}连续导热中的温度场与熵产生&\href{course-page:8:122}{8-119}\\[2pt]
\hspace{1em}习题与解答&\href{course-page:8:125}{8-122}\\[2pt]
\textbf{第四讲　相变、界面与外场平衡}&\href{course-page:8:132}{8-129}\\[2pt]
\hspace{1em}选择状态变量与守恒的能量&\href{course-page:8:132}{8-129}\\[2pt]
\hspace{1em}表面张力系统&\href{course-page:8:137}{8-134}\\[2pt]
\hspace{1em}相平衡与汽化热&\href{course-page:8:142}{8-139}\\[2pt]
\hspace{1em}相空间分布与外场中的粒子&\href{course-page:8:147}{8-144}\\[2pt]
\hspace{1em}界面、相变与外场的综合模型&\href{course-page:8:155}{8-152}\\[2pt]
\hspace{1em}磁性工质的熵与绝热降温&\href{course-page:8:165}{8-162}\\[2pt]
\hspace{1em}习题与解答&\href{course-page:8:167}{8-164}\\[2pt]
\end{longtable}
\section*{第九编\quad 近代物理进阶}
\begin{longtable}{@{}p{.84\textwidth}r@{}}
\textbf{第一讲　事件、四维量与运动光源}&\href{course-page:9:5}{9-1}\\[2pt]
\hspace{1em}从测量程序确定事件&\href{course-page:9:5}{9-1}\\[2pt]
\hspace{1em}四维矢量与固有量&\href{course-page:9:9}{9-5}\\[2pt]
\hspace{1em}速度、光行差与方向变换&\href{course-page:9:12}{9-8}\\[2pt]
\hspace{1em}四速度、四动量与四力&\href{course-page:9:14}{9-10}\\[2pt]
\hspace{1em}运动光源的角分布&\href{course-page:9:16}{9-12}\\[2pt]
\hspace{1em}坐标测量中的事件选择&\href{course-page:9:17}{9-13}\\[2pt]
\hspace{1em}四维量的几何和角分布的测度&\href{course-page:9:19}{9-15}\\[2pt]
\hspace{1em}有限装置中的同时事件与雷达测距&\href{course-page:9:21}{9-17}\\[2pt]
\hspace{1em}光锥、波面和速度空间&\href{course-page:9:25}{9-21}\\[2pt]
\hspace{1em}角分布积分与四维微分关系&\href{course-page:9:27}{9-23}\\[2pt]
\hspace{1em}相对速率的标量表达与辐射的局部能量&\href{course-page:9:29}{9-25}\\[2pt]
\hspace{1em}习题与解答&\href{course-page:9:32}{9-28}\\[2pt]
\textbf{第二讲　时空过程、相对论火箭与延迟成像}&\href{course-page:9:46}{9-42}\\[2pt]
\hspace{1em}坐标原点与仿射变换&\href{course-page:9:46}{9-42}\\[2pt]
\hspace{1em}共线变换与快度&\href{course-page:9:49}{9-45}\\[2pt]
\hspace{1em}从守恒定律推导火箭方程&\href{course-page:9:50}{9-46}\\[2pt]
\hspace{1em}延迟时间与视运动&\href{course-page:9:54}{9-50}\\[2pt]
\hspace{1em}连续转向与加速参考系&\href{course-page:9:57}{9-53}\\[2pt]
\hspace{1em}推进方程的能量和质量含义&\href{course-page:9:58}{9-54}\\[2pt]
\hspace{1em}延迟成像的投影几何&\href{course-page:9:59}{9-55}\\[2pt]
\hspace{1em}往返过程中的钟和坐标网&\href{course-page:9:61}{9-57}\\[2pt]
\hspace{1em}火箭的有限点火与燃料分配&\href{course-page:9:64}{9-60}\\[2pt]
\hspace{1em}近场接收与远场视速度&\href{course-page:9:67}{9-63}\\[2pt]
\hspace{1em}有限加速装置的距离与信号&\href{course-page:9:69}{9-65}\\[2pt]
\hspace{1em}同一实验室时刻怎样读取加速杆两端的钟&\href{course-page:9:72}{9-68}\\[2pt]
\hspace{1em}习题与解答&\href{course-page:9:74}{9-70}\\[2pt]
\textbf{第三讲　相对论动力学与电磁场}&\href{course-page:9:87}{9-83}\\[2pt]
\hspace{1em}相位与波四矢量&\href{course-page:9:87}{9-83}\\[2pt]
\hspace{1em}总四动量、质心系与阈值&\href{course-page:9:88}{9-84}\\[2pt]
\hspace{1em}四力、加速度与瞬时静止系&\href{course-page:9:90}{9-86}\\[2pt]
\hspace{1em}四电流与电荷守恒&\href{course-page:9:92}{9-88}\\[2pt]
\hspace{1em}四势、规范与场方程&\href{course-page:9:93}{9-89}\\[2pt]
\hspace{1em}电磁场的完整变换&\href{course-page:9:94}{9-90}\\[2pt]
\hspace{1em}守恒方程中的未知量与参考系&\href{course-page:9:96}{9-92}\\[2pt]
\hspace{1em}四势推导中的全部横向分量&\href{course-page:9:97}{9-93}\\[2pt]
\hspace{1em}麦克斯韦方程与电荷连续性&\href{course-page:9:98}{9-94}\\[2pt]
\hspace{1em}从三力到轨迹&\href{course-page:9:100}{9-96}\\[2pt]
\hspace{1em}电荷、电流与场的同一变换&\href{course-page:9:102}{9-98}\\[2pt]
\hspace{1em}散射与衰变中的不变质量&\href{course-page:9:105}{9-101}\\[2pt]
\hspace{1em}弹性碰撞的能量转移与交叉场轨迹&\href{course-page:9:108}{9-104}\\[2pt]
\hspace{1em}两体末态的能谱与连续光压推进&\href{course-page:9:111}{9-107}\\[2pt]
\hspace{1em}习题与解答&\href{course-page:9:114}{9-110}\\[2pt]
\textbf{第四讲　光量子、热辐射与半经典量子模型}&\href{course-page:9:126}{9-122}\\[2pt]
\hspace{1em}光电效应与表面势垒&\href{course-page:9:126}{9-122}\\[2pt]
\hspace{1em}腔辐射的模态与能量分布&\href{course-page:9:128}{9-124}\\[2pt]
\hspace{1em}从经典配分到普朗克谱&\href{course-page:9:129}{9-125}\\[2pt]
\hspace{1em}玻尔氢样原子与相对论圆轨道&\href{course-page:9:132}{9-128}\\[2pt]
\hspace{1em}作用量量子化与转折边界&\href{course-page:9:133}{9-129}\\[2pt]
\hspace{1em}从波长关系到束缚态方程&\href{course-page:9:137}{9-133}\\[2pt]
\hspace{1em}WKB相位、振幅和量子偏移&\href{course-page:9:138}{9-134}\\[2pt]
\hspace{1em}零点能和位置—动量展宽&\href{course-page:9:139}{9-135}\\[2pt]
\hspace{1em}黑体谱中经典近似的边界&\href{course-page:9:140}{9-136}\\[2pt]
\hspace{1em}光子统计与辐射能量尺度&\href{course-page:9:141}{9-137}\\[2pt]
\hspace{1em}能级间隔、周期与边界相位&\href{course-page:9:145}{9-141}\\[2pt]
\hspace{1em}有限束缚边界与势阱的时间改变&\href{course-page:9:150}{9-146}\\[2pt]
\hspace{1em}自由原子的吸收共振与反冲&\href{course-page:9:154}{9-150}\\[2pt]
\hspace{1em}习题与解答&\href{course-page:9:155}{9-151}\\[2pt]
\end{longtable}
\section*{第十编\quad 光学进阶}
\begin{longtable}{@{}p{.84\textwidth}r@{}}
\textbf{第一讲　近轴成像、渐变介质与光程}&\href{course-page:10:5}{10-1}\\[2pt]
\hspace{1em}符号与单球面光焦度&\href{course-page:10:5}{10-1}\\[2pt]
\hspace{1em}约化光线矩阵与复合系统&\href{course-page:10:7}{10-3}\\[2pt]
\hspace{1em}费马原理与力学类比&\href{course-page:10:9}{10-5}\\[2pt]
\hspace{1em}复振幅与光程相位&\href{course-page:10:12}{10-8}\\[2pt]
\hspace{1em}费马原理的变分推导&\href{course-page:10:13}{10-9}\\[2pt]
\hspace{1em}主平面如何连接实际光线&\href{course-page:10:15}{10-11}\\[2pt]
\hspace{1em}渐变介质力学类比的参数选择&\href{course-page:10:15}{10-11}\\[2pt]
\hspace{1em}成像矩阵与横向放大率&\href{course-page:10:17}{10-13}\\[2pt]
\hspace{1em}渐变介质的传像与精确光线&\href{course-page:10:20}{10-16}\\[2pt]
\hspace{1em}光程、相位与短波近似&\href{course-page:10:23}{10-19}\\[2pt]
\hspace{1em}有限孔径、光瞳与可成像的视场&\href{course-page:10:24}{10-20}\\[2pt]
\hspace{1em}矩阵成像与高斯光束的连接&\href{course-page:10:28}{10-24}\\[2pt]
\hspace{1em}习题与解答&\href{course-page:10:31}{10-27}\\[2pt]
\textbf{第二讲　双束干涉、薄膜与多束相干叠加}&\href{course-page:10:46}{10-42}\\[2pt]
\hspace{1em}杨氏干涉与光程差&\href{course-page:10:46}{10-42}\\[2pt]
\hspace{1em}薄膜的斜入射程差与反射相位&\href{course-page:10:51}{10-47}\\[2pt]
\hspace{1em}多束几何级数与分辨能力&\href{course-page:10:53}{10-49}\\[2pt]
\hspace{1em}相干叠加的时间平均&\href{course-page:10:57}{10-53}\\[2pt]
\hspace{1em}薄膜程差的同波面比较&\href{course-page:10:59}{10-55}\\[2pt]
\hspace{1em}焦散相移的波动来由&\href{course-page:10:59}{10-55}\\[2pt]
\hspace{1em}有限束数与多次反射腔的比较&\href{course-page:10:61}{10-57}\\[2pt]
\hspace{1em}有限带宽、有限光源与曝光时间&\href{course-page:10:62}{10-58}\\[2pt]
\hspace{1em}多次反射与干涉测量&\href{course-page:10:64}{10-60}\\[2pt]
\hspace{1em}腔内损耗、耦合匹配与相位测量&\href{course-page:10:68}{10-64}\\[2pt]
\hspace{1em}谱线形状、相位扩散与色散包络&\href{course-page:10:71}{10-67}\\[2pt]
\hspace{1em}习题与解答&\href{course-page:10:74}{10-70}\\[2pt]
\textbf{第三讲　衍射积分、孔径与光栅分辨率}&\href{course-page:10:85}{10-81}\\[2pt]
\hspace{1em}电磁波、复数平均与光强&\href{course-page:10:85}{10-81}\\[2pt]
\hspace{1em}从惠更斯积分到远场傅里叶积分&\href{course-page:10:85}{10-81}\\[2pt]
\hspace{1em}单缝的积分与相量圆弧&\href{course-page:10:86}{10-82}\\[2pt]
\hspace{1em}双缝包络与有限宽光栅&\href{course-page:10:89}{10-85}\\[2pt]
\hspace{1em}分辨判据与孔径宽度&\href{course-page:10:90}{10-86}\\[2pt]
\hspace{1em}闪耀反射光栅&\href{course-page:10:92}{10-88}\\[2pt]
\hspace{1em}衍射积分中的尺度与归一化&\href{course-page:10:93}{10-89}\\[2pt]
\hspace{1em}傅里叶积分中的能量守恒&\href{course-page:10:94}{10-90}\\[2pt]
\hspace{1em}点像、物像与分辨判据&\href{course-page:10:95}{10-91}\\[2pt]
\hspace{1em}近场边缘与孔径的空间频率&\href{course-page:10:97}{10-93}\\[2pt]
\hspace{1em}连续孔径、二维光栅与采样&\href{course-page:10:100}{10-96}\\[2pt]
\hspace{1em}孔径自相关、成像传递与边缘加权&\href{course-page:10:104}{10-100}\\[2pt]
\hspace{1em}孔面相位误差与相干峰的损失&\href{course-page:10:108}{10-104}\\[2pt]
\hspace{1em}习题与解答&\href{course-page:10:111}{10-107}\\[2pt]
\textbf{第四讲　偏振、晶体与界面电磁波}&\href{course-page:10:125}{10-121}\\[2pt]
\hspace{1em}Jones矢量与偏振椭圆&\href{course-page:10:125}{10-121}\\[2pt]
\hspace{1em}自然光、部分偏振与检偏&\href{course-page:10:127}{10-123}\\[2pt]
\hspace{1em}单轴晶体中的普通光与非普通光&\href{course-page:10:128}{10-124}\\[2pt]
\hspace{1em}波片及偏振态变换&\href{course-page:10:129}{10-125}\\[2pt]
\hspace{1em}界面边界条件与菲涅耳公式&\href{course-page:10:132}{10-128}\\[2pt]
\hspace{1em}能流、互易与布儒斯特角&\href{course-page:10:132}{10-128}\\[2pt]
\hspace{1em}全反射相位与倏逝波&\href{course-page:10:134}{10-130}\\[2pt]
\hspace{1em}偏振椭圆的代数结构&\href{course-page:10:136}{10-132}\\[2pt]
\hspace{1em}从Jones矢量到相干矩阵&\href{course-page:10:137}{10-133}\\[2pt]
\hspace{1em}单轴晶体色散式的推导&\href{course-page:10:138}{10-134}\\[2pt]
\hspace{1em}菲涅耳能量关系的逐项证明&\href{course-page:10:139}{10-135}\\[2pt]
\hspace{1em}全反射的解析支与相位连续性&\href{course-page:10:139}{10-135}\\[2pt]
\hspace{1em}从测量强度还原偏振态&\href{course-page:10:141}{10-137}\\[2pt]
\hspace{1em}晶体中的波矢与能流&\href{course-page:10:143}{10-139}\\[2pt]
\hspace{1em}全反射延迟与边界附近能流&\href{course-page:10:145}{10-141}\\[2pt]
\hspace{1em}偏振元件的损耗与色散&\href{course-page:10:147}{10-143}\\[2pt]
\hspace{1em}闭合偏振路径中的干涉相位&\href{course-page:10:150}{10-146}\\[2pt]
\hspace{1em}习题与解答&\href{course-page:10:152}{10-148}\\[2pt]
\end{longtable}

\chapter*{主题与模型索引}\addcontentsline{toc}{chapter}{主题与模型索引}
同一主题下列出相关节题及其起始页。需要比较不同模型时，可由这些入口查阅相应推导、例题与习题。
\begingroup\small\linespread{1.1}\selectfont
\setlength{\parskip}{2pt}
\ctexset{section={format=\normalsize\bfseries,beforeskip=12pt,afterskip=6pt}}
\section*{复数、相量与频率响应}
\noindent 第一编·从二维几何到复数运算\quad\href{course-page:1:6}{1-1}\par
\noindent 第一编·同频振动与相量\quad\href{course-page:1:11}{1-6}\par
\noindent 第十编·电磁波、复数平均与光强\quad\href{course-page:10:85}{10-81}\par
\noindent 第十编·单缝的积分与相量圆弧\quad\href{course-page:10:86}{10-82}\par
\section*{曲率、自然坐标与相对运动}
\noindent 第一编·曲线的切向、法向与曲率\quad\href{course-page:1:76}{1-71}\par
\noindent 第二编·曲率与自然坐标\quad\href{course-page:2:92}{2-90}\par
\noindent 第二编·动坐标与相对运动\quad\href{course-page:2:113}{2-111}\par
\noindent 第三编·梁的弯曲、曲率与端部条件\quad\href{course-page:3:329}{3-326}\par
\noindent 第八编·相对运动、碰撞与自由程\quad\href{course-page:8:15}{8-12}\par
\section*{约束、自由度与广义坐标}
\noindent 第一编·从约束选坐标：抛物线轨道上的小珠\quad\href{course-page:1:81}{1-76}\par
\noindent 第二编·绳杆约束与刚体关联\quad\href{course-page:2:47}{2-45}\par
\noindent 第二编·自由度与坐标选择\quad\href{course-page:2:210}{2-208}\par
\noindent 第三编·能量积分、坐标加速度与约束反力\quad\href{course-page:3:4}{3-1}\par
\noindent 第三编·多自由度、循环坐标与整体反冲\quad\href{course-page:3:21}{3-18}\par
\noindent 第四编·约束运动中的降阶与终止条件\quad\href{course-page:4:75}{4-73}\par
\noindent 第四编·无约束极值与约束极值\quad\href{course-page:4:116}{4-114}\par
\noindent 第五编·局部展开与约束下的稳定性\quad\href{course-page:5:49}{5-47}\par
\section*{虚功、平衡与稳定性}
\noindent 第一编·二次型与平衡点稳定性\quad\href{course-page:1:137}{1-132}\par
\noindent 第二编·平衡点与分界轨线\quad\href{course-page:2:156}{2-154}\par
\noindent 第二编·虚功与静力平衡\quad\href{course-page:2:218}{2-216}\par
\noindent 第三编·约束极值、虚功与二阶变化\quad\href{course-page:3:63}{3-60}\par
\noindent 第三编·平衡构形、二次能量与固有频率\quad\href{course-page:3:135}{3-132}\par
\noindent 第五编·局部展开与约束下的稳定性\quad\href{course-page:5:49}{5-47}\par
\noindent 第八编·相变、界面与外场平衡\quad\href{course-page:8:132}{8-129}\par
\noindent 第八编·相平衡与汽化热\quad\href{course-page:8:142}{8-139}\par
\section*{守恒量与系统边界}
\noindent 第二编·守恒关系与参数平面的允许区域\quad\href{course-page:2:179}{2-177}\par
\noindent 第三编·多自由度、循环坐标与整体反冲\quad\href{course-page:3:21}{3-18}\par
\noindent 第三编·碰撞守恒量、冲量与有效质量\quad\href{course-page:3:180}{3-177}\par
\noindent 第四编·变换中的守恒关系与收敛条件\quad\href{course-page:4:199}{4-197}\par
\noindent 第八编·选择状态变量与守恒的能量\quad\href{course-page:8:132}{8-129}\par
\noindent 第九编·从守恒定律推导火箭方程\quad\href{course-page:9:50}{9-46}\par
\noindent 第九编·四电流与电荷守恒\quad\href{course-page:9:92}{9-88}\par
\noindent 第十编·傅里叶积分中的能量守恒\quad\href{course-page:10:94}{10-90}\par
\section*{相图、势垒与临界运动}
\noindent 第二编·相图的应用\quad\href{course-page:2:147}{2-145}\par
\noindent 第九编·光电效应与表面势垒\quad\href{course-page:9:126}{9-122}\par
\section*{简正模、矩阵与网络谱}
\noindent 第一编·矩阵、线性变换与方程组\quad\href{course-page:1:127}{1-122}\par
\noindent 第一编·非正交坐标与度量矩阵\quad\href{course-page:1:132}{1-127}\par
\noindent 第三编·简正模问题\quad\href{course-page:3:135}{3-132}\par
\noindent 第三编·振型分解、简正坐标与初始条件\quad\href{course-page:3:149}{3-146}\par
\noindent 第七编·有限矩形网络的边界谱与端口电阻\quad\href{course-page:7:82}{7-79}\par
\noindent 第七编·习题2.6：横纵电阻不同的六节点网\quad\href{course-page:7:90}{7-87}\par
\noindent 第十编·约化光线矩阵与复合系统\quad\href{course-page:10:7}{10-3}\par
\noindent 第十编·成像矩阵与横向放大率\quad\href{course-page:10:17}{10-13}\par
\section*{共振、阻尼与能量交换}
\noindent 第一编·受迫阻尼振动：从瞬态到稳态\quad\href{course-page:1:16}{1-11}\par
\noindent 第一编·阻尼模型的完整比较\quad\href{course-page:1:19}{1-14}\par
\noindent 第二编·阻尼及综合应用\quad\href{course-page:2:165}{2-163}\par
\noindent 第三编·拍振、失谐与外力响应\quad\href{course-page:3:159}{3-156}\par
\noindent 第九编·自由原子的吸收共振与反冲\quad\href{course-page:9:154}{9-150}\par
\section*{碰撞、冲量与变质量运动}
\noindent 第二编·连续受力与碰撞散射\quad\href{course-page:2:20}{2-18}\par
\noindent 第三编·碰撞守恒量、冲量与有效质量\quad\href{course-page:3:180}{3-177}\par
\noindent 第三编·多约束、摩擦与连续碰撞\quad\href{course-page:3:196}{3-193}\par
\noindent 第八编·相对运动、碰撞与自由程\quad\href{course-page:8:15}{8-12}\par
\noindent 第八编·碰撞抽样与运动边界的进一步应用\quad\href{course-page:8:27}{8-24}\par
\noindent 第九编·从守恒定律推导火箭方程\quad\href{course-page:9:50}{9-46}\par
\noindent 第九编·火箭的有限点火与燃料分配\quad\href{course-page:9:64}{9-60}\par
\section*{刚体、滚动与进动}
\noindent 第二编·绳杆约束与刚体关联\quad\href{course-page:2:47}{2-45}\par
\noindent 第二编·接触与滚动关联\quad\href{course-page:2:55}{2-53}\par
\noindent 第三编·纯滚动系统的动能与接触条件\quad\href{course-page:3:16}{3-13}\par
\noindent 第三编·滑动摩擦与滚动状态的转换\quad\href{course-page:3:26}{3-23}\par
\section*{微积分、渐近展开与误差}
\noindent 第一编·曲线坐标中的微分算子（选读）\quad\href{course-page:1:101}{1-96}\par
\noindent 第一编·矩阵指数与线性微分方程（选读）\quad\href{course-page:1:155}{1-150}\par
\noindent 第四编·导数、微分与物理过程\quad\href{course-page:4:3}{4-1}\par
\noindent 第四编·泰勒展开与保留阶数\quad\href{course-page:4:12}{4-10}\par
\noindent 第五编·场的微分与曲线坐标\quad\href{course-page:5:13}{5-11}\par
\noindent 第五编·边界误差怎样传入内部区域\quad\href{course-page:5:70}{5-68}\par
\noindent 第六编·局部场方程与微分运算\quad\href{course-page:6:6}{6-3}\par
\noindent 第七编·非均匀泄漏带的幂律与局部近似\quad\href{course-page:7:127}{7-124}\par
\noindent 第九编·角分布积分与四维微分关系\quad\href{course-page:9:27}{9-23}\par
\section*{变分法与最值}
\noindent 第三编·曲线约束与最速下降\quad\href{course-page:3:71}{3-68}\par
\noindent 第四编·二次变分与曲线坐标的几何\quad\href{course-page:4:141}{4-139}\par
\noindent 第十编·费马原理的变分推导\quad\href{course-page:10:13}{10-9}\par
\section*{静电边值、唯一性与镜像}
\noindent 第五编·源、场与电势\quad\href{course-page:5:3}{5-1}\par
\noindent 第五编·边界条件与唯一性\quad\href{course-page:5:46}{5-44}\par
\noindent 第六编·电势、磁矢势与规范\quad\href{course-page:6:7}{6-4}\par
\noindent 第七编·导电与静电的边值类比\quad\href{course-page:7:17}{7-14}\par
\noindent 第七编·唯一性与等势节点的归并\quad\href{course-page:7:47}{7-44}\par
\section*{电介质、极化与磁介质}
\noindent 第五编·极化强度、电位移与束缚电荷\quad\href{course-page:5:132}{5-130}\par
\section*{电磁场能量、动量与应力}
\noindent 第五编·由场能到导体表面的电压力\quad\href{course-page:5:24}{5-22}\par
\noindent 第六编·磁能的形变导数与电磁应力\quad\href{course-page:6:158}{6-155}\par
\section*{电磁感应、磁链与自感}
\noindent 第六编·感应电压的测量路径与滑动接触\quad\href{course-page:6:26}{6-23}\par
\noindent 第六编·磁通、磁链与电感的定义\quad\href{course-page:6:48}{6-45}\par
\section*{粒子轨道与导体运动}
\noindent 第一编·从约束选坐标：抛物线轨道上的小珠\quad\href{course-page:1:81}{1-76}\par
\noindent 第三编·由初态确定轨道的形状与方向\quad\href{course-page:3:356}{3-353}\par
\noindent 第三编·轨道微扰、进动与潮汐\quad\href{course-page:3:377}{3-374}\par
\noindent 第六编·带电粒子的能量与正则动量\quad\href{course-page:6:10}{6-7}\par
\noindent 第六编·磁场边界、转动与受约束轨道\quad\href{course-page:6:160}{6-157}\par
\noindent 第七编·横档与轨道不等阻时的衰减尺度\quad\href{course-page:7:101}{7-98}\par
\noindent 第九编·玻尔氢样原子与相对论圆轨道\quad\href{course-page:9:132}{9-128}\par
\section*{电阻网络的对称性与互易}
\noindent 第一编·立体角、通量与对称性\quad\href{course-page:1:89}{1-84}\par
\noindent 第一编·事件终点的灵敏度与参数优化\quad\href{course-page:1:203}{1-198}\par
\noindent 第六编·互感、互易与耦合强度\quad\href{course-page:6:56}{6-53}\par
\noindent 第六编·几何对称性与互感的分解\quad\href{course-page:6:100}{6-97}\par
\noindent 第七编·对称性之外的桥式网络\quad\href{course-page:7:53}{7-50}\par
\noindent 第七编·互易定理与多端口\quad\href{course-page:7:67}{7-64}\par
\noindent 第七编·习题2.4：电桥失衡后的电流与灵敏度\quad\href{course-page:7:88}{7-85}\par
\noindent 第七编·习题2.5：互易网络接入有限负载\quad\href{course-page:7:89}{7-86}\par
\noindent 第七编·习题2.7：各支路有效电阻的加权和\quad\href{course-page:7:91}{7-88}\par
\noindent 第十编·能流、互易与布儒斯特角\quad\href{course-page:10:132}{10-128}\par
\section*{差分、无限网络与分形}
\noindent 第七编·轮形网络与无限递推\quad\href{course-page:7:75}{7-72}\par
\noindent 第七编·有限梯形与端部条件\quad\href{course-page:7:93}{7-90}\par
\noindent 第七编·习题3.5：两种横档交替的半无限梯形\quad\href{course-page:7:135}{7-132}\par
\noindent 第七编·习题3.6：由有限梯形的测量反求参数\quad\href{course-page:7:136}{7-133}\par
\section*{暂态电路与初值}
\noindent 第五编·边界几何造成的场增强与角模衰减\quad\href{course-page:5:66}{5-64}\par
\noindent 第七编·横档与轨道不等阻时的衰减尺度\quad\href{course-page:7:101}{7-98}\par
\noindent 第七编·串联RLC的二阶暂态\quad\href{course-page:7:156}{7-153}\par
\noindent 第七编·习题4.6：由衰减峰值反求串联电路参数\quad\href{course-page:7:181}{7-178}\par
\noindent 第七编·习题4.7：初始充电选择衰减模态\quad\href{course-page:7:183}{7-180}\par
\section*{分子分布与输运}
\noindent 第一编·连续分布的积分：圆环与圆盘\quad\href{course-page:1:89}{1-84}\par
\noindent 第二编·角分布、立体角与靶面密度\quad\href{course-page:2:12}{2-10}\par
\noindent 第三编·质量输运、动量与能量收支\quad\href{course-page:3:101}{3-98}\par
\noindent 第三编·弯曲波的色散与能量输运\quad\href{course-page:3:333}{3-330}\par
\noindent 第六编·感生电场的空间分布\quad\href{course-page:6:12}{6-9}\par
\noindent 第八编·从状态概率到速度分布\quad\href{course-page:8:4}{8-1}\par
\noindent 第八编·速率、方向与抽样方式\quad\href{course-page:8:9}{8-6}\par
\noindent 第九编·运动光源的角分布\quad\href{course-page:9:16}{9-12}\par
\section*{热力学状态函数与偏导数}
\noindent 第四编·偏导数、全微分与复合函数\quad\href{course-page:4:112}{4-110}\par
\noindent 第八编·泄流、辐射与热力学关系\quad\href{course-page:8:48}{8-45}\par
\noindent 第八编·热力学势与偏导关系\quad\href{course-page:8:57}{8-54}\par
\section*{熵、热机与可用功}
\noindent 第八编·热机、卡诺循环与效率上界\quad\href{course-page:8:90}{8-87}\par
\noindent 第八编·熵的定义与不可逆过程\quad\href{course-page:8:93}{8-90}\par
\section*{相变、相平衡与外场}
\noindent 第八编·相平衡与汽化热\quad\href{course-page:8:142}{8-139}\par
\noindent 第八编·相空间分布与外场中的粒子\quad\href{course-page:8:147}{8-144}\par
\section*{相对论、时空与快度}
\noindent 第三编·固有加速度与相对论推进\quad\href{course-page:3:112}{3-109}\par
\noindent 第九编·共线变换与快度\quad\href{course-page:9:49}{9-45}\par
\noindent 第九编·玻尔氢样原子与相对论圆轨道\quad\href{course-page:9:132}{9-128}\par
\section*{光子、物质波与量子束缚}
\noindent 第五编·极化强度、电位移与束缚电荷\quad\href{course-page:5:132}{5-130}\par
\noindent 第八编·光子气体的压强与辐射\quad\href{course-page:8:55}{8-52}\par
\noindent 第九编·作用量量子化与转折边界\quad\href{course-page:9:133}{9-129}\par
\noindent 第九编·从波长关系到束缚态方程\quad\href{course-page:9:137}{9-133}\par
\section*{几何光学与成像}
\noindent 第十编·成像矩阵与横向放大率\quad\href{course-page:10:17}{10-13}\par
\noindent 第十编·有限孔径、光瞳与可成像的视场\quad\href{course-page:10:24}{10-20}\par
\section*{干涉、衍射与相干}
\noindent 第十编·杨氏干涉与光程差\quad\href{course-page:10:46}{10-42}\par
\noindent 第十编·相干叠加的时间平均\quad\href{course-page:10:57}{10-53}\par
\noindent 第十编·衍射积分、孔径与光栅分辨率\quad\href{course-page:10:85}{10-81}\par
\noindent 第十编·衍射积分中的尺度与归一化\quad\href{course-page:10:93}{10-89}\par
\section*{偏振、晶体与能流}
\noindent 第十编·Jones矢量与偏振椭圆\quad\href{course-page:10:125}{10-121}\par
\noindent 第十编·自然光、部分偏振与检偏\quad\href{course-page:10:127}{10-123}\par
\noindent 第十编·单轴晶体中的普通光与非普通光\quad\href{course-page:10:128}{10-124}\par
\noindent 第十编·全反射相位与倏逝波\quad\href{course-page:10:134}{10-130}\par
\noindent 第十编·晶体中的波矢与能流\quad\href{course-page:10:143}{10-139}\par
\section*{数值计算与参数反演}
\noindent 第一编·简谐振子的数值误差\quad\href{course-page:1:184}{1-179}\par
\noindent 第一编·参数反演中的可识别性\quad\href{course-page:1:199}{1-194}\par
\noindent 第七编·习题3.6：由有限梯形的测量反求参数\quad\href{course-page:7:136}{7-133}\par
\noindent 第七编·习题4.6：由衰减峰值反求串联电路参数\quad\href{course-page:7:181}{7-178}\par

\endgroup
\end{document}
